\documentclass[11pt]{article}

\usepackage[margin=0.85in]{geometry}
\usepackage[T1]{fontenc}
\usepackage[utf8]{inputenc}
\usepackage[strings]{underscore}
\usepackage{booktabs}
\usepackage{longtable}
\usepackage{array}
\usepackage{graphicx}
\usepackage{hyperref}
\usepackage{enumitem}
\usepackage{xcolor}
\usepackage{listings}
\usepackage[capitalise]{cleveref}

\hypersetup{colorlinks=true,linkcolor=blue,citecolor=blue,urlcolor=blue}

\lstdefinestyle{pyblock}{
  basicstyle=\ttfamily\small,
  breaklines=true,
  columns=fullflexible,
  keywordstyle=\color{blue!60!black},
  commentstyle=\color{gray},
  stringstyle=\color{green!40!black},
  frame=single,
  language=Python,
}
\lstset{style=pyblock}

\newcolumntype{L}[1]{>{\raggedright\arraybackslash}p{#1}}

\title{\textbf{ACI-CO Repository Architecture}}
\author{Colombian Actuarial Climate Index (ACI-CO) pipeline}
\date{\today}

\begin{document}
\sloppy
\maketitle

\begin{abstract}
This document is a LaTeX companion to \texttt{ARCHITECTURE.md} and
\texttt{CLAUDE.md} at the repository root. It describes how the ACI-CO
codebase is organized: the scientific target it computes, the actual data
flow between scripts (as opposed to the README's idealized version), where
the repository's documentation is authoritative versus stale, and how to
read already-processed outputs from a notebook (e.g.\ Google Colab) using
the \texttt{aci_lib} data-access library introduced alongside this document.
It is a description of the codebase as it stands, not a specification --
where the code and the docs disagree, the code wins.
\end{abstract}

\tableofcontents

%======================================================================
\section{What this repository is}
%======================================================================

This is a research pipeline that builds the \textbf{Colombian Actuarial
Climate Index (ACI-CO)} from ERA5(-Land) reanalysis data (1961--2024): a
Colombian adaptation of the Actuaries Climate Index (ACI) framework,
extended with ENSO calibration, a bootstrap uncertainty band, multi-dataset
robustness checks, and validation against Colombia's national disaster
registry (UNGRD).

It is an active research codebase, not a packaged library: there is no test
suite, several parallel/experimental implementations of the same pipeline
stage coexist deliberately, and a non-trivial share of the \emph{working}
pipeline was never committed to git (\Cref{sec:git-gap}).

\paragraph{Authoritative methodology source.} The methodology (how T90,
T10, Rx5day, CDD and WP are computed; how the composite, ENSO decomposition,
and bootstrap bands work) is described in
\texttt{articles/\allowbreak aci\_co\_submission/\allowbreak article1.tex}, \S3--\S7 -- \emph{not}
the older root-level \texttt{articles/article1.tex} (a superseded
LOWESS-based draft), and not the empty stubs in \texttt{docs/methodology.md}
/ \texttt{docs/results.md} / \texttt{docs/zones\_description.md}.

%======================================================================
\section{The scientific target: ACI-CO}
%======================================================================

The ACI-CO composite is the unweighted mean of five signed, standardized
components (divided by $K=5$; an earlier draft divided by 6 for a
placeholder sea-level component, since dropped for lack of tide-gauge
coverage):

\begin{longtable}{@{}L{2.2cm}L{6.2cm}L{1.5cm}L{4.5cm}@{}}
\toprule
\textbf{Component} & \textbf{Meaning} & \textbf{Sign} & \textbf{Threshold} \\
\midrule
\endhead
T90 & Fraction of days/month with temperature above the 1961--1990 baseline 90th percentile & $+$ & Empirical P90 per grid cell $\times$ calendar month \\
T10 & Fraction of days/month with temperature below the baseline 10th percentile & $-$ & Empirical P10 \\
Rx5day & Max 5-day precipitation total per month, $z$-scored vs.\ 1961--1990 climatology & $+$ & $z$-score vs.\ monthly climatology \\
CDD & Annual max consecutive dry days ($<$1\,mm), interpolated onto the monthly axis & $+$ & --- \\
WP & Wind-power ($\propto U^3$) exceedance & $+$ & Empirical P90 (current); replaces an earlier parametric $\mu + 1.28\sigma$ threshold, discarded because WP is right-skewed \\
\bottomrule
\end{longtable}

\subsection{Extensions beyond the base ACI framework}

\begin{itemize}[leftmargin=*]
\item \textbf{ENSO-neutral decomposition} (\S4): asymmetric ARIMAX
  regression per component/region on Ni\~no-3.4 ($E^+=\max(E,0)$,
  $E^-=\min(E,0)$) plus a linear trend and AR($p$) errors, separating
  secular warming from ENSO-driven variability.
\item \textbf{Bootstrap uncertainty bands} (\S4.4): block-resample
  ($L=12$\,mo) the 1961--1990 baseline composite into $B=500$ (national) /
  $B=200$ (regional) null-model replicates, testing whether the observed
  trend exceeds baseline-period variability. Replaced an earlier LOWESS
  residual-envelope approach.
\item \textbf{Multi-dataset robustness} (\S5): ERA5-Land vs.\ ERA5 standard
  vs.\ CHIRPS.
\item \textbf{UNGRD validation} (\S7): negative-binomial regression of
  disaster counts (floods, landslides, windthrow, vegetation fires --
  49{,}734 events, 1998--2022) on ACI-CO components.
\end{itemize}

\subsection{Regions actually computed}
\label{sec:aci-co-target}

Despite the paper conceptually discussing 5 natural regions and 32
departments, the implementation aggregates \textbf{4 units}: national
(area-weighted, all of Colombia) and 3 departments chosen for insurance
relevance -- \textbf{Antioquia}, \textbf{Cundinamarca-Bogot\'a}, and
\textbf{Valle del Cauca}. Shapefiles also exist for Bogot\'a, Medell\'in,
Cali, San Andr\'es/Providencia, Pac\'ifico, and Amazonas, but the
corresponding \texttt{data/processed/anomalias\_\{bogota,medellin,\allowbreak san\_andres\_providencia\}/}
directories are empty -- never populated.

%======================================================================
\section{Repository layout}
%======================================================================

\begin{longtable}{@{}L{3.2cm}L{10.8cm}@{}}
\toprule
\textbf{Path} & \textbf{Contents} \\
\midrule
\endhead
\texttt{src/scripts/} & The original OPT pipeline: download, percentile, and anomaly-calculation scripts (\texttt{calcular\_*}, \texttt{ecmwf\_descarga.py}, \texttt{unir\_archivos.py}, \texttt{graficas.py}, \dots). Spanish-language docstrings dominate here. \\
\texttt{src/forecast\_scripts/} & Two independent forecasting families (AutoTS ensembles vs.\ SARIMAX/ETS + ENSO) plus \texttt{common/}, the shared modules extracted from both (\Cref{sec:forecast-common}). English-language, newer code. \\
\texttt{src/dashboard/} & \texttt{app.py} -- a read-only Streamlit dashboard that reads Stage 3/5 CSVs directly; not a pipeline stage. \\
\texttt{src/utils/} & Shared helpers (\texttt{get\_available\_years}, \texttt{get\_cached\_shapefile}), consolidated from near-identical copies previously duplicated across the \texttt{calcular\_*} scripts. \\
\texttt{src/aci\_lib/} & Read-only data-access library for notebooks (\Cref{sec:aci-lib}) -- added alongside this document, not part of the original OPT/forecast pipeline. \\
\texttt{data/raw/}, \texttt{data/processed/} & Pipeline inputs/outputs, entirely gitignored except for \texttt{.gitkeep} placeholders (\Cref{sec:git-gap}). \\
\texttt{articles/} & Three papers at different stages: the current submission (\texttt{aci\_co\_submission/}), a superseded LOWESS draft (\texttt{article1.tex}), and a separate forecasting report (\texttt{AnnualICAForecast\_Report.tex}). \\
\texttt{docs/} & Mostly empty stubs (\texttt{methodology.md}, \texttt{results.md}, \texttt{zones\_description.md}); the one substantive file is \texttt{Diccionario.md}, a $\sim$2400-line data-source vetting document. \\
\texttt{tests/} & Placeholder only (\texttt{test.txt}); there is no test suite or CI. \\
\bottomrule
\end{longtable}

%======================================================================
\section{Actual data flow}
%======================================================================

The README describes an idealized 8-step flow; several of the scripts it
names do not exist (\Cref{sec:discrepancies}). The flow below is
reconstructed from what the scripts actually read and write.

\begin{longtable}{@{}L{2.2cm}L{5.2cm}L{6.6cm}@{}}
\toprule
\textbf{Stage} & \textbf{Script(s)} & \textbf{Output} \\
\midrule
\endhead
0. Download & \texttt{ecmwf\_descarga.py} (canonical), \texttt{sst\_descarga.py}, \texttt{get\_aux\_data.py}, IDEAM station data via \texttt{pyreadr} & \texttt{data/raw/era5/*.grib}, SST GRIB, \texttt{data/raw/auxiliary/}, \texttt{data/raw/ideam/*.rds} \\
1. Merge / resample & \texttt{unir\_archivos.py}; independently, the untracked \texttt{run\_aci\_colombia.py} re-merges raw GRIB into a from-scratch reimplementation & \texttt{era5\_daily\_combined\_\{tmp,\allowbreak rain,\allowbreak wind\}.nc}; \texttt{data/processed/aci\_daily/}, \texttt{aci\_colombia\_1961\_2024.csv} \\
2. Percentiles (1961--1990 baseline) & \texttt{calcular\_percentil\_\{temperatura,\allowbreak lluvia,\allowbreak viento\}.py} & \texttt{era5\_*\_percentil.nc}, \texttt{percentiles*.csv} \\
3. Anomalies, per region & \texttt{calcular\_anomalias\_\{temperatura,\allowbreak lluvia,\allowbreak viento\}.py}, orchestrated by \texttt{calcular\_anomalias\_regiones.py} over \texttt{data/shapefiles/*.shp} & \texttt{data/processed/anomalias\_<region>/anomalies\_\{temperature,\allowbreak precipitation,\allowbreak drought,\allowbreak wind\}\_combined.csv} (populated: antioquia, cali, colombia, cundinamarca\_bogota, valle\_cauca; empty: bogota, medellin, san\_andres\_providencia) \\
4. Daily variants & \texttt{produce\_daily\_anomalies.py}, \texttt{append\_2025\_2026.py} (untracked) & \texttt{data/processed/daily\_anomalies/*.xlsx} \\
5. ENSO / SST index & \texttt{anomalias\_sst\_daily.py}, \texttt{compare\_enso\_flags.py} & \texttt{sst\_index\_colombia\_pacific.csv}, \texttt{enso\_flag\_comparison.csv}, validated against NOAA CPC \texttt{oni.ascii.txt} \\
6. Downscaling & \texttt{downscale\_era5\_temperature.py}, \texttt{downscale\_era5\_tp.py} (RandomForest, DEM/NDVI/IDEAM) & \texttt{data/processed/downscaled/}, validated in \texttt{data/processed/era5\_ideam\_comparison/} \\
7. Visualization / export & \texttt{graficas.py}, \texttt{generate\_correlation\_matrices.py} & \texttt{articles/graficas/<region>/*.png}, xlsx exports, correlation heatmaps \\
8. Forecasting & Two independent families -- see \Cref{sec:forecast-common} & \texttt{articles/graficas/forecast\_*/} \\
\bottomrule
\end{longtable}

\texttt{data/indices/} and \texttt{data/zones/} exist but are effectively
empty -- no script currently writes to those exact paths.

\paragraph{Raw ERA5 \texttt{.grib} backup.} \texttt{data/raw/era5/*.grib}
is gitignored (\Cref{sec:git-gap}) and not reproduced anywhere in version
control. A backup copy of these raw downloads is kept on Google Drive:
\url{https://drive.google.com/drive/folders/17SdbsxVJYnqjPzIlS4FFZ6MowfMcHfqW}.

\subsection{Two parallel ``correctness check'' implementations}

There are effectively two independent implementations of the core ACI
algorithm in this repository:

\begin{itemize}[leftmargin=*]
\item \textbf{The OPT pipeline} (\texttt{src/scripts/calcular\_*}, tracked
  in git, the README's documented flow) -- the original project pipeline,
  per-region, percentile/anomaly based.
\item \textbf{\texttt{run\_aci\_colombia.py}} (untracked) -- a from-scratch
  reimplementation of the ACI-Python reference algorithm directly on raw
  GRIB, used to cross-validate the OPT pipeline's output.
  \texttt{compare\_aci\_colombia.py} and \texttt{compare\_repos.py} (also
  untracked) diff the two and document the algorithmic differences.
\end{itemize}

Neither is ``the old version'' of the other -- they are independent
implementations kept deliberately in sync for cross-validation, and this
comparison grounds the paper's ``validation against x-ACI'' discussion
(\S8).

\subsection{Forecasting: two families plus shared modules}
\label{sec:forecast-common}

Two independent forecasting families exist side by side and are \emph{not}
a superset/subset of each other:

\begin{itemize}[leftmargin=*]
\item \texttt{forecast\_ica\_\{monthly,daily\}.py} -- AutoTS ensembles.
\item The \texttt{ETS(X)\_*} family -- statsmodels SARIMAX/ETS with an
  ENSO exogenous regressor.
\end{itemize}

\texttt{src/forecast\_scripts/common/} holds logic that was byte-identical
or near-identical across these scripts before consolidation:

\begin{longtable}{@{}L{3.0cm}L{5.7cm}L{5.1cm}@{}}
\toprule
\textbf{Module} & \textbf{Extracted from} & \textbf{Used by} \\
\midrule
\endhead
\texttt{data\_loading.py} & Byte-identical in \texttt{forecast\_ica\_monthly.py} and \texttt{ETS\_ica\_forecast.py} & Both, via thin same-named wrapper methods \\
\texttt{climate\_index.py} & Near-identical in all three monthly forecasters & All three \\
\texttt{stationarity.py} & Near-identical in \texttt{ETS\_ica\_forecast.py} / \texttt{ETSX\_ica\_forecast.py} & Both \\
\texttt{diagnostics.py} (\texttt{Model\-Error\-Diagnostics}) & Moved unchanged from the old \texttt{error\_diagnostics.py} & Not currently called by anything \\
\bottomrule
\end{longtable}

Deliberately \emph{not} merged (structurally parallel, not literal
duplicates): the \texttt{ONIDataHandler}\allowbreak/\texttt{ONIDataHandlerDaily}
ENSO-fetching classes, the SARIMAX+Lasso fitting methods, the AutoTS setup
(same shape, different hyperparameters), and all
\texttt{visualize\_forecasts}-style plotting methods (4 separate
implementations, each with different overlays).

%======================================================================
\section{Documentation status}
%======================================================================

\begin{longtable}{@{}L{4.6cm}L{9.4cm}@{}}
\toprule
\textbf{Document} & \textbf{Status} \\
\midrule
\endhead
\texttt{docs/methodology.md}, \texttt{docs/results.md}, \texttt{docs/zones\_description.md} & Empty stubs. Methodology lives in \texttt{articles/aci\_co\_submission/article1.tex} instead. \\
\texttt{docs/explicacion\_excel.md} & One truncated sentence, no body. \\
\texttt{docs/data\_dictionary.md} & Early, incomplete fragment (HUMBOLDT-only) of \texttt{docs/Diccionario.md}. \\
\texttt{docs/Diccionario.md} & The real, complete data-source vetting document ($\sim$2400 lines, 20 candidate sources evaluated). Canonical over \texttt{data\_dictionary.md}. \\
\texttt{README.md} & Describes an idealized 8-step flow; two referenced scripts do not exist. \\
\bottomrule
\end{longtable}

%======================================================================
\section{Discrepancies between README/docstrings and reality}
\label{sec:discrepancies}
%======================================================================

\begin{itemize}[leftmargin=*]
\item README step 1 says to run \texttt{descargar\_datos.py} -- it does not
  exist; \texttt{ecmwf\_descarga.py} is the actual entry point.
\item README step 6 says to run \texttt{sealevel.py} -- does not exist.
  \texttt{sealevel2.py} exists but does not do sea-level analysis; it is a
  Mann-Kendall trend + Q-Q normality check on precipitation-like series with
  leftover Jupyter cell markers, suggesting it was pasted from a notebook
  and misnamed. Sea level is not part of the current ACI-CO methodology.
\item \texttt{src/scripts/earthub\_descarga.py} has a whole-file
  indentation error and cannot run as-is, on top of embedding a live-looking
  credential (\Cref{sec:secrets}).
\item \texttt{src/scripts/display.py} is not a plotting module -- it is a
  Greece-precipitation download script (also embeds the same credential)
  that produced a 562\,MB \texttt{.nc} file at the repo root as a side
  effect.
\item Stale internal docstrings/comments: \texttt{era5datasets.py} says
  \texttt{"""compare\_land\_only.py"""}; \texttt{get\_aux\_data.py}'s header
  says \texttt{get\_auxiliary\_data.py}. Trust the filename/git path.
\item \texttt{graficas\_altaresolucion.py} is a 4-line dead stub.
\item \texttt{ETSX\_seasonal\_motif\_forecast.py} (the ``broken two-stage
  Lasso$\to$ETS'' approach \texttt{ETSX\_daily\_ica\_forecast.py}'s
  docstring says it replaced) has been deleted -- confirmed unused before
  removal.
\item \texttt{example\_grid\_search\_era5.py} imports \texttt{ONIDataHandler}
  from \texttt{ETSX\_ica\_forecast.py} expecting methods that do not exist
  on that class -- a pre-existing break, unrelated to the forecast\_scripts
  modularization.
\item \texttt{articles/AnnualICAForecast\_Report.tex} references
  \texttt{scripts/daily\_forecast\_to\_annual.py} and
  \texttt{scripts/ar1\_vs\_ets\_final.py} -- neither exists anywhere in the
  repository. Treat that report's pipeline as documented-but-unimplemented.
\end{itemize}

%======================================================================
\section{The git tracking gap}
\label{sec:git-gap}
%======================================================================

A significant portion of the \emph{current, more-correct} pipeline exists
only in the working tree and has never been \texttt{git add}ed:
\texttt{run\_aci\_colombia.py}, \texttt{compare\_aci\_colombia.py},
\texttt{compare\_repos.py}, \texttt{create\_region\_shapefiles.py},
\texttt{probe\_2025.py}, \texttt{produce\_daily\_anomalies.py},
\texttt{append\_2025\_2026.py}, \texttt{reproduce.py},
\texttt{plot\_repo\_anomalies.py}, and most of \texttt{articles/} (the
\texttt{aci\_co\_submission/} folder, \texttt{produce\_paper\_plots.py},
\texttt{build\_validation\_outputs.py}, etc.).

Conversely, things that \emph{are} committed include generated junk
(\texttt{presentation.aux/log/out}, \texttt{texput.log},
\texttt{output.log}, \texttt{comparison.png}, \texttt{storm\_daniel.png},
\texttt{.vs/}, \texttt{proyecto-indices/}, root-level \texttt{.nc}/\texttt{.csv}
data files) and a \texttt{.gitignore} that was itself accidentally
overwritten with a PowerShell heredoc command instead of its output.

\textbf{Untracked here often means ``uncommitted work in progress,'' not
``disposable.''} Do not delete untracked files without checking which side
of this split they are on.

\subsection{Known secrets in the repository}
\label{sec:secrets}

Two Earth Data Hub PAT tokens are hardcoded and \textbf{committed} in
\texttt{src/scripts/display.py} (line 13) and
\texttt{src/scripts/earthub\_descarga.py} (line 15) -- the same token,
duplicated. A CDS API key also appears in
\texttt{GRID\_SEARCH\_AND\_ERA5\_GUIDE.md} (line 44). These require
rotation/history-scrubbing independent of any code change. New download
scripts should follow the pattern in \texttt{ecmwf\_descarga.py} /
\texttt{sst\_descarga.py}: \texttt{cdsapi.Client()} reads credentials from
the user's local \texttt{\textasciitilde/.cdsapirc}, nothing embedded in
source.

%======================================================================
\section{The \texttt{aci\_lib} data-access library}
\label{sec:aci-lib}
%======================================================================

\texttt{src/aci\_lib/} is a small, dependency-light Python package that
turns \texttt{data/processed/} into a browsable, loadable catalog, so that
already-computed pipeline outputs can be pulled into a notebook (locally or
in Google Colab) without knowing the directory's internal layout by heart.

\subsection{Why it exists}

\texttt{data/processed/} mixes single tabular files, single gridded
NetCDF files, and long runs of per-year or per-year-month NetCDF files that
really form one logical time series (e.g.\ 768 files under
\texttt{anomalias\_colombia/anomalies\_temperature\_\allowbreak 1961\_1.nc}
through \texttt{\_2024\_12.nc}). \texttt{aci\_lib} groups the latter
automatically and exposes a uniform three-function API:

\begin{itemize}[leftmargin=*]
\item \texttt{list\_datasets()} -- a DataFrame of every dataset name, kind,
  file count, and (where known) a description sourced from this
  document's data-flow tables.
\item \texttt{describe(name)} -- the underlying file list and kind for one
  dataset.
\item \texttt{load(name)} -- returns a \texttt{pandas.DataFrame} for
  CSV/XLSX datasets, an \texttt{xarray.Dataset} for single or multi-file
  NetCDF datasets (opened lazily via \texttt{xarray.open\_mfdataset}), or a
  file path for images/raw text files.
\end{itemize}

\subsection{Important caveat: data is not in git}

\texttt{data/processed/} is entirely gitignored (\Cref{sec:git-gap}) --
only a \texttt{.gitkeep} placeholder is tracked. A plain \texttt{git clone}
of this repository, including in Colab, brings the \emph{code} in
\texttt{src/aci\_lib/} but \emph{not} the data it reads. The data must be
supplied separately: mount it from Google Drive, upload and unzip an
archive of \texttt{data/processed/}, or copy it in via any other means
available in the notebook environment.

\subsection{Quickstart in Google Colab}

\begin{lstlisting}
# 1. Get the code (data/processed/ is gitignored, so this clone
#    does NOT include the data -- only src/aci_lib/).
!git clone https://github.com/<org>/aca_indice_climatico_opt-main.git
import sys
sys.path.insert(0, "/content/aca_indice_climatico_opt-main/src")

# 2. Supply the data. Pick ONE of:
#  (a) mount Drive, if data/processed/ was uploaded there beforehand:
from google.colab import drive
drive.mount("/content/drive")
DATA_DIR = "/content/drive/MyDrive/aca_data/processed"

#  (b) upload + unzip a local copy:
# from google.colab import files
# files.upload()  # e.g. processed.zip
# !unzip -q processed.zip -d /content/aci_data
# DATA_DIR = "/content/aci_data/processed"

import aci_lib
aci_lib.set_data_dir(DATA_DIR)

# 3. Browse and load.
catalog = aci_lib.list_datasets()
catalog.head(20)

national_composite = aci_lib.load(
    "anomalias_colombia/anomalies_temperature_combined"
)  # -> pandas.DataFrame

gridded_series = aci_lib.load(
    "anomalias_colombia/anomalies_temperature"
)  # -> xarray.Dataset, lazily concatenated from ~768 monthly files
\end{lstlisting}

Run locally, inside a checkout of this repository, \texttt{set\_data\_dir()}
is unnecessary: \texttt{aci\_lib.get\_data\_dir()} walks up from
\texttt{src/aci\_lib/} to find the repository root (marked by
\texttt{CLAUDE.md}) and locates \texttt{data/processed/} on its own.

\subsection{Scope}

\texttt{aci\_lib} currently covers everything under \texttt{data/processed/}
only -- it does not reach into \texttt{articles/*/graficas/\allowbreak paper\_figures/}
(the paper's final figures/tables) or \texttt{data/raw/}. Both are
reasonable follow-up extensions if a notebook workflow needs them, but were
kept out of scope to keep the catalog's grouping logic simple and its
descriptions trustworthy.

%======================================================================
\section{Benchmarking the pipeline in Colab}
\label{sec:benchmark}
%======================================================================

Unlike \texttt{aci\_lib} (read-only over already-computed
\texttt{data/processed/}), this capability runs the actual pipeline
\emph{stages} -- Stage 1 merge/resample, Stage 2 baseline percentiles,
Stage 3 anomalies -- against raw ERA5 \texttt{.grib}, to see how long they
take on whatever compute a given Colab session provides.

\subsection{Why no script changes were needed}

\texttt{unir\_archivos.py}, \texttt{calcular\_percentil\_temperatura.py},
and \texttt{calcular\_anomalias\_temperatura.py} already expose plain,
parameterized functions (\texttt{process\_yearly\_data\_tmp},
\texttt{calcular\_percentiles}, \texttt{procesar\_anomalias\_temperatura},
\dots) behind a thin \texttt{\_\_main\_\_} CLI wrapper that just supplies
default repo-relative paths. \textbf{They were already importable as a
library} -- calling them directly from a notebook with explicit paths
works without modification. This was verified by importing
\texttt{unir\_archivos} and timing
\texttt{process\_yearly\_precipitation\_data} against a real local
\texttt{.grib} file before writing the Colab notebook described below.

\subsection{\texttt{src/drive\_sync.py}}

A new, small helper module (not part of \texttt{aci\_lib}) for pulling data
out of the shared Google Drive folder (\Cref{sec:git-gap}) into this
repo's \texttt{data/} layout, for use inside Colab. That folder
(\texttt{My Drive/2. Datos}, confirmed from the Drive UI -- it is a
shortcut sitting directly in \texttt{My Drive} root; ``Indice Climatico
Actuarial'', seen in Drive's breadcrumb/location panel, is the owner's
canonical path, not where the shortcut lives) is shared with specific
named collaborators, not ``anyone with the link'' -- confirmed from its
``Who has access'' panel, which is also why anonymous fetching (tested
with \texttt{gdown}) returns HTTP 401. So \texttt{drive\_sync.sync(...)}
instead copies (or symlinks) from an already-mounted
\texttt{google.colab.drive}, which authenticates as whoever is logged into
that Colab session.

The full structure of \texttt{2. Datos} is now confirmed (via
\texttt{notebooks/explore\_drive\_structure.ipynb} in the
\texttt{aca\_aci\_collab} repo):
\texttt{shapefiles/} is flat (48 files, matches \texttt{data/shapefiles/});
\texttt{era5/completos/} is the flat raw archive (195 files,
\texttt{era5\_\{rain,tmp,wind\}\_<year>.grib} -- exactly what
\texttt{data/raw/era5/} expects); \texttt{era5/union/} and
\texttt{era5/Combined/<date>/} hold Stage 1 merged-daily outputs
(differently named, or date-versioned, respectively -- not a direct copy
target); \texttt{era5/percentiles\_nc/<date>/} holds dated Stage 2
baseline-percentile snapshots (names \emph{do} match
\texttt{data/processed/} exactly, but picking ``current'' means picking a
date); \texttt{salidas\_colombia/} and
\texttt{salidas\_cundinamarca\_bogota/} are each a single \texttt{.zip},
not a flat folder; and \texttt{datos\_indice/} is a dated personal archive
(\texttt{2024-12/}, \texttt{2025-06/}, \texttt{2025-12/}) of
\texttt{anomalias\_<region>} CSVs that doesn't map onto any one local
directory.

\texttt{DEFAULT\_MAPPING} accordingly only covers the three
unambiguous cases -- \texttt{shapefiles}, \texttt{era5/completos}
(nested-path keys work: \texttt{sync()} just \texttt{os.path.join}s them),
and the two \texttt{salidas\_*} zips, which \texttt{sync()} detects and
extracts automatically rather than copying the \texttt{.zip} itself. The
dated/renamed subfolders (\texttt{union/}, \texttt{Combined/},
\texttt{percentiles\_nc/}, \texttt{datos\_indice/}) are deliberately left
out -- picking a date or renaming a file is a judgment call the module
does not make for you; pass an explicit \texttt{mapping=} for those.

\subsection{\texttt{aca\_opt\_benchmark.ipynb}}

Lives in the \texttt{aca\_aci\_collab} companion repo's
\texttt{notebooks/}, alongside \texttt{aci\_lib\_quickstart.ipynb}, but
clones \emph{this} pipeline repo rather than the companion repo, since
that's what has the scripts to benchmark. It fetches one year of raw
temperature \texttt{.grib} (not the whole $\sim$13\,GB archive) and times
Stages 1--3 against just that year, then extrapolates linearly to a full
1961--2024 (64-year) run for Stages 1 and 3. Stage 2 (baseline percentiles)
is timed against the same single year purely as a computational-shape
proxy -- the real pipeline computes it once over the fixed 30-year
1961--1990 baseline, not once per year, so its extrapolated figure isn't
meaningful; multiply the single-year Stage 2 time by roughly 30 instead.

\subsection{Verified, then partially reopened: Colab vs.\ local
reproducibility}
\label{sec:reproducibility}

\textbf{1-year smoke test (2026-08-10, year 1985, single process):}
Stage 1 took 31.8s in Colab vs.\ 22.1s locally, Stage 2 11.7s vs.\ 8.8s,
Stage 3 30.8s vs.\ 26.6s -- different hardware, same order of magnitude.
The computed \texttt{t\_90}/\texttt{t\_10} anomaly values matched
\textbf{bit-for-bit} between the two runs ($-0.436035$ / $-0.737441$ for
December 1985, identical \texttt{NaN} pattern for the other 11 months).
An unpinned Colab environment (\texttt{xarray==2026.7.0} vs.\ this repo's
\texttt{xarray==2024.11.0}) silently dropped the
\texttt{count\_hot}/\texttt{count\_cold} columns in this run; pinning
every relevant package to \texttt{requirements.txt}'s exact versions
(\texttt{xarray==2024.11.0}, \texttt{cfgrib==0.9.14.1},
\texttt{rioxarray==0.18.2}, \texttt{geopandas==1.0.1},
\texttt{netCDF4==1.7.2}, \texttt{eccodes==2.39.1}) appeared to resolve it
at the time.

\textbf{Full run (2026-08-10, 1961--2024, 30-year baseline,
\texttt{use\_multiprocessing=True}):} with the same pinned versions, the
\texttt{count\_hot}/\texttt{count\_cold} columns were \emph{still}
missing, and \texttt{t\_90}/\texttt{t\_10} no longer matched exactly
(max abs.\ diff $\approx 0.02$, mean $\approx 0.001$, against the official
\texttt{salidas\_colombia} reference, 768/768 rows aligned). This
contradicts the smoke test's pinning-fixes-it conclusion above -- either
the pin silently didn't take in that session, or version pinning was
never the real cause. \textbf{Leading hypothesis, not yet confirmed:} the
smoke test ran with \texttt{use\_multiprocessing=False}; the full run
(like the real pipeline's own default) used
\texttt{multiprocessing.Pool}. \texttt{calcular\_anomalias\_temperatura.py}
keeps a module-level \texttt{\_percentiles\_cache} dict wrapping lazily-
loaded \texttt{xr.open\_dataset()} handles -- a known-fragile pattern
under fork-based multiprocessing, where file handles duplicated across
\texttt{fork()} can behave inconsistently. Re-running Stage 3 with
\texttt{use\_multiprocessing=False} at full scale would confirm or rule
this out; not yet done. Treat both the missing columns and the
$t\_90$/$t\_10$ drift as \textbf{open, unresolved} until then -- do not
trust downstream (composite index, multi-region) results computed with
\texttt{use\_multiprocessing=True} without re-checking this.

\subsection{\texttt{aca\_opt\_full\_replication.ipynb}}

Scales \texttt{aca\_opt\_benchmark.ipynb} up from a 1-year smoke test to
the real thing: fetches all \texttt{era5\_tmp\_<year>.grib} files
(1961--2024) rather than one, runs Stage 1/2 over the actual 30-year
1961--1990 baseline (not a 1-year proxy), Stage 3 over the full record,
and then diffs the resulting \texttt{anomalies\_temperature\_combined.csv}
against the official \texttt{salidas\_colombia} output already on Drive
(fetched into a separate \texttt{\_official\_reference/} path so it can't
be overwritten by the fresh computation) -- column-by-column max/mean
absolute difference, not just a visual spot-check. Expected runtime is
45--90 minutes (network fetch of $\sim$5.5\,GB plus compute); not
something to run casually, but the strongest available evidence that
\texttt{aca\_indice\_climatico\_opt} genuinely works as an importable
library against Drive-sourced data end to end, not just for one
hand-picked year.

\subsection{\texttt{aca\_opt\_multi\_region\_monthly.ipynb}}

Extends the replication from temperature-only to all four monthly
components \texttt{calcular\_anomalias\_regiones.py}'s existing
\texttt{procesar\_anomalias\_region()} already orchestrates per region --
temperature (T90/T10), wind power (WP), precipitation (Rx5day), and
drought (CDD) -- across the four regions this repo documents as actually
populated (\Cref{sec:aci-co-target}: national, Antioquia,
Cundinamarca-Bogot\'a, Valle del Cauca). No new pipeline logic; this wires
already-importable functions to Drive-sourced data across three raw
variables (temperature, precipitation, wind) instead of one. Expected
runtime is 2.5--3.5 hours for the default 4 regions -- substantially
longer than the temperature-only replication, since it triples the raw
data fetched and runs three anomaly calculations per region.

\textbf{Two implementation details surfaced while wiring this up, both
pre-existing repo quirks, not new bugs:}
\begin{itemize}[leftmargin=*]
\item \texttt{unir\_archivos.py}'s three \texttt{process\_yearly\_*\_data}
  functions have \textbf{inconsistent parameter order}:
  \texttt{process\_yearly\_data\_tmp} and
  \texttt{process\_yearly\_precipitation\_data} take
  \texttt{(file\_path, year, variable, output\_dir)}, but
  \texttt{process\_yearly\_wind\_data} takes
  \texttt{(file\_path, year, output\_dir, variable)} -- \texttt{output\_dir}
  and \texttt{variable} swapped. A generic wrapper assuming one order
  would silently misdirect wind's output; the notebook uses small
  per-variable adapters instead of "fixing" the scripts' order.
\item \texttt{calcular\_percentil\_lluvia.py}'s \texttt{calcular\_percentiles()}
  saves \texttt{era5\_lluvia\_percentil.nc} (singular), but
  \texttt{calcular\_anomalias\_lluvia.py} reads
  \texttt{era5\_lluvias\_percentil.nc} (plural) -- confirmed as a real,
  pre-existing mismatch: this repo's own \texttt{data/processed/}
  contains \emph{both} filenames side by side, evidently from someone
  working around it manually rather than fixing the source. The notebook
  copies the percentile output under both names rather than silently
  hiding the inconsistency; fixing the actual name mismatch in the two
  scripts would be a reasonable, separate follow-up.
\end{itemize}

\textbf{A third issue found this way was an actual bug, not a quirk, and
was fixed in the source (not worked around):}
\texttt{calcular\_anomalias\_lluvia.py}'s \texttt{process\_year()} selected
a year's grib file with \texttt{[f for f in files if str(year) in f and
f.endswith(".grib")][0]} -- no variable-name filter, unlike
\texttt{calcular\_anomalias\_temperatura.py} and
\texttt{calcular\_anomalias\_viento.py}, which both filter further
(\texttt{"tmp" in f} / \texttt{"wind" in f}). As long as
\texttt{data/raw/era5/} held only rain grib this never mattered; once this
notebook fetched all three variables into one directory, it started
picking temperature or wind files for most years, failing with
\texttt{"No variable named 'tp'"}. Caught live during an actual Colab run
(2026-08-11); the missing \texttt{"rain" in f} filter was added directly
to the script (matching the existing pattern in the other two), verified
against real data before pushing. The already-published
\texttt{salidas\_colombia} reference was unaffected (768/768 rows present,
confirmed by row count) -- it was evidently computed back when
\texttt{data/raw/era5/} held only rain grib, before temperature/wind were
added.

\subsection{Performance: decode raw grib once per year, not once per
region}

The first version of this notebook called
\texttt{procesar\_anomalias\_region()} once per region, and each call
independently re-decoded the same raw grib from scratch -- 4 regions
$\times$ 3 variables $\times$ 64 years = 768 redundant grib-decode passes,
identical until the final shapefile clip -- understandably frustrating
given how long a run was taking, since multiprocessing across years does
not fix redundant work, and GPU/TPU runtimes provide \emph{zero} benefit
here (nothing in this stack -- \texttt{xarray}/\texttt{cfgrib}/
\texttt{pandas} -- is CUDA-aware).

Temperature and wind both expose a clean split --
\texttt{load\_annual\_grid\_data} / \texttt{load\_annual\_grid\_data\_safe}
take an \emph{optional} \texttt{shapefile\_path}, decoding unclipped data
when it's \texttt{None} -- so the notebook now decodes each year once and
clips that one decode to all 4 regions via plain \texttt{rioxarray}
(\texttt{.rio.write\_crs(...).rio.clip(...)}), calling
\texttt{calcular\_anomalias}/\texttt{calcular\_anomalias\_viento} per
region-month against the shared decode. Multiprocessing still parallelizes
over years (unchanged from before), so this combines with, not replaces,
the existing 7-worker parallelism. \textbf{Verified bit-for-bit identical}
against the original per-region code path (year 1985, Colombia, both
temperature and wind) before shipping -- same values, roughly a 4x
reduction in redundant decode work.

\subsection{Stage 1 had no parallelism at all -- an oversight, not a
hardware limit}

Separately from the region-level redundancy above,
\texttt{aca\_opt\_multi\_region\_monthly.ipynb}'s (and
\texttt{aca\_opt\_full\_replication.ipynb}'s) Stage 1 originally processed
its 30 baseline years in a plain sequential \texttt{for} loop --
\texttt{unir\_archivos.py}'s three \texttt{process\_yearly\_*\_data}
functions have no cross-year dependency (each year writes a uniquely
named intermediate file, unlike drought's CDD in
\Cref{sec:cdd-race}), so this was safely, trivially parallelizable and
simply wasn't. This alone accounted for roughly 30 of the multi-region
notebook's total minutes. Fixed by parallelizing over years with the same
\texttt{multiprocessing.Pool} pattern used elsewhere -- named top-level
task functions per variable, since \texttt{Pool} pickles tasks to send to
worker processes (even under Colab's \texttt{fork} start method) and
lambdas aren't picklable. Verified with real \texttt{Pool} execution (not
just reasoning about it) against 6 real years of local data: 34.2s
parallel vs.\ $\sim$132s sequential, all 6 output files correct.

\subsection{Precipitation/drought: a real correctness hazard, not just a
performance gap}
\label{sec:cdd-race}

Precipitation is \emph{not} sped up the same decode-once way temperature
and wind are, and drought could not safely be either, for a reason more
serious than "not yet done":
\texttt{calcular\_anomalias\_lluvia.py}'s \texttt{calcular\_interpolacion()}
(the CDD/drought calculation) reads the \textbf{previous year's} saved
\texttt{datos\_maximos\_\{year-1\}.nc} to compute the current year's value,
and writes its own for the next year -- a sequential, year-over-year
dependency. \texttt{procesar\_anomalias\_lluvia()}'s own default is
\texttt{use\_multiprocessing=True}, parallelizing over years via
\texttt{multiprocessing.Pool} -- which does not guarantee execution order.
Any year whose predecessor has not yet been written when it runs silently
falls back to a cruder approximation (\texttt{except FileNotFoundError}
in \texttt{calcular\_interpolacion}) instead of the proper interpolation.
\textbf{This is a pre-existing risk in this script's own default
behavior, not something introduced by this notebook} -- and it is
plausible, not yet confirmed, that some official reference output was
affected, since \texttt{use\_multiprocessing=True} is the default a plain
call would use.

(\texttt{load\_grid\_data}'s \texttt{shapefile} parameter was still made
properly optional -- its docstring already said "if a shapefile is
provided, clips the data," but the implementation clipped unconditionally
with no \texttt{None} guard. That part was a real, narrow, worth-fixing
gap; the CDD ordering issue is a different, larger one.)

Given that, \texttt{aca\_opt\_multi\_region\_monthly.ipynb} parallelizes
precipitation/drought across \textbf{regions} instead of years: each
region gets its own worker process running
\texttt{procesar\_anomalias\_lluvia(..., use\_multiprocessing=False)} --
sequential and therefore correct within that region -- with up to 4
regions running concurrently. Real parallelism (up to 4x, one region per
core) without touching the ordering question above. Verified by calling
the worker function directly (not through \texttt{Pool}, since Windows'
\texttt{spawn} multiprocessing model does not reliably reflect Colab's
Linux \texttt{fork} model for local testing) against two regions --
distinct, plausible output for both, no cross-region collisions.

This leaves the CDD ordering issue itself unresolved -- a legitimate
follow-up would be making \texttt{calcular\_interpolacion} not depend on
sequential file I/O at all (e.g.\ computing each year's raw dry-day count
independently and doing the cross-year blend as a separate, explicitly
ordered pass), which is a real redesign, not something to bundle into a
performance-motivated notebook change.

This changes performance only for temperature/wind, and the region- vs.\
year-level parallelization choice only for precipitation/drought -- the
open multiprocessing-correctness question from \Cref{sec:reproducibility}
(the \texttt{count\_hot}/\texttt{count\_cold} and \texttt{t\_90}/\texttt{t\_10}
drift, specific to temperature) is unaffected by any of this and remains
open regardless of which version of this notebook is run.

\subsection{Rendering \texttt{article1.tex} inline}

A cell in \texttt{aca\_opt\_multi\_region\_monthly.ipynb} compiles
\texttt{articles/aci\_co\_submission/article1.tex} -- the authoritative
methodology paper (\Cref{sec:aci-co-target}), \emph{not} the superseded
root-level \texttt{articles/article1.tex} -- and embeds the resulting PDF
inline via a base64 data URI, rather than just offering it for download.
This only works because that file (and its figures under
\texttt{graficas/paper\_figures/}) is actually tracked in git -- confirmed
before building this, since \texttt{ARCHITECTURE.md}'s git-tracking-gap
section originally listed all of \texttt{articles/aci\_co\_submission/} as
untracked; that changed once the pending local commits were pushed
(\Cref{sec:git-gap}). \texttt{article1.tex} uses a hand-written
\texttt{thebibliography} block, not \texttt{natbib}+\texttt{bibtex}+a
separate \texttt{.bib} file, so two \texttt{pdflatex} passes (for
cross-references/TOC) compile it with no \texttt{bibtex}/\texttt{biber}
step needed -- confirmed by compiling it locally (53 pages, no undefined
references, no LaTeX errors) before writing the notebook cell.

\subsection{Daily vs.\ monthly: a cross-validation, not just a resolution
comparison}

\texttt{produce\_daily\_anomalies.py} (daily, DOY-bootstrap thresholds,
following \texttt{run\_aci\_colombia.py}'s algorithm) and
\texttt{calcular\_anomalias\_temperatura.py} (monthly, calendar-month
thresholds, the OPT pipeline used everywhere else in this notebook) are
\textbf{two independently-implemented pipelines} kept in sync specifically
to cross-check each other (\Cref{sec:git-gap} on
\texttt{run\_aci\_colombia.py}'s original purpose), not two resolutions of
the same computation. A new notebook section compares them properly:

\begin{itemize}[leftmargin=*]
\item Reuses \texttt{run\_aci\_colombia.py}'s per-year
  \emph{preprocessing} functions only (\texttt{preprocess\_temperature\_-
  year} etc.) -- not its full two-phase pipeline -- against the same raw
  grib already fetched for the OPT pipeline, avoiding a second download.
  Verified against real local data (1961--1962): both years' temperature,
  wind, and rain preprocessing completed correctly ($\sim$88s/year for all
  three variables combined).
\item Compares \textbf{unstandardized} exceedance fractions -- the OPT
  pipeline's \texttt{count\_hot} column against the daily pipeline's
  \texttt{T90} resampled to a monthly mean -- not the \emph{standardized}
  \texttt{t\_90} column used everywhere else in this notebook, which is
  not on the same scale as either.
\item Scoped to Colombia, temperature only, the 1961--1990 baseline (the
  DOY thresholds need it) plus a configurable \texttt{COMPARISON\_YEARS}
  (default: 2020--2024) -- widening it costs preprocessing time linearly.
\item \texttt{produce\_daily\_anomalies.py}'s own \texttt{doy\_percentile}
  computes each of the 366 day-of-year thresholds as a separate
  \texttt{.isel()}/\texttt{.quantile()} xarray call -- measured at
  \textbf{150s for one of the four calls needed, on just a 2-year test
  subset} (would be far worse against the real 30-year baseline). A
  functionally-identical numpy rewrite (same day-of-year windowing logic,
  \texttt{np.percentile} instead of 366 separate xarray operations) measured
  \textbf{1.6s} on the same input -- a 92x speedup, with a $1.83\times
  10^{-5}$ max difference (floating-point-level, from a different default
  interpolation method, not a real disagreement). The notebook uses this
  faster version, defined locally in the cell -- \texttt{produce\_daily\_-
  anomalies.py} itself is unchanged, since this is a notebook-side
  convenience, not a pipeline correctness fix. The full chain (fast
  thresholds $\to$ \texttt{region\_mask} $\to$
  \texttt{compute\_t90\_t10\_daily} $\to$ \texttt{spatial\_mean}) was
  verified end-to-end against real local data before shipping: 6.9s for
  all four thresholds, sensible T90 values in the expected $[0,1]$ range.
\end{itemize}

\subsection{Drive-backed result caching, raw-grib access, and loading
outputs back}

Every restart this session re-ran Stage 1/2 (the 1961--1990 baseline
merge/percentiles) from scratch -- $\sim$20--30 minutes each time, for
outputs identical across runs regardless of which regions or
\texttt{COMPARISON\_YEARS} Stage 3 uses. Three additions address this and
two adjacent, explicitly-requested needs:

\begin{itemize}[leftmargin=*]
\item \textbf{Caching} (\texttt{drive\_sync.cache\_upload}/
  \texttt{cache\_restore}/\texttt{cache\_complete}, verified locally --
  empty-cache detection, upload, complete-cache detection, restore into a
  fresh directory, partial-cache correctly reported as incomplete):
  deliberately \textbf{file-list-based, not whole-folder}, and
  deliberately scoped to \emph{only} the 8 small, stable Stage 1/2
  baseline files -- \emph{not} Stage 3's per-region output, where this
  session's still-open correctness questions live
  (\Cref{sec:reproducibility}, \Cref{sec:cdd-race}). Caching those
  silently would risk reusing something still being debugged. Defaults to
  the user's own \texttt{My Drive} space, not the shared \texttt{2. Datos}
  team folder -- writing generated cache files into shared storage by
  default, without teammates expecting them, isn't something to do
  unprompted.
\item \textbf{Raw grib access} (\texttt{load\_raw\_grib()}): a thin
  wrapper around \texttt{xr.open\_dataset(..., engine="cfgrib")} plus
  optional shapefile clipping, for exploration outside the fixed
  pipeline -- verified against real data, both unclipped (71$\times$68
  full extent) and clipped to a small region (Valle del Cauca, correctly
  reduced to 8$\times$6).
\item \textbf{Loading outputs for further analysis}: rather than
  reimplementing a browse/load layer, the notebook clones
  \texttt{aca\_aci\_collab} and points \texttt{aci\_lib}
  (\Cref{sec:aci-lib}) directly at the session's own
  \texttt{data/processed/} -- what gets loaded is what this run just
  computed, not a separately-fetched copy.
\end{itemize}

%======================================================================
\section{Cleanup priorities (standing list)}
%======================================================================

Ordered roughly by urgency/impact; see \texttt{ARCHITECTURE.md} for the
canonical, living version of this list.

\begin{enumerate}[leftmargin=*]
\item \textbf{Secrets}: rotate the Earth Data Hub PAT tokens and the CDS
  key (\Cref{sec:secrets}); scrub from git history if the two committed
  copies need to be fully removed.
\item \textbf{Commit the real pipeline}: the untracked files listed in
  \Cref{sec:git-gap}.
\item \textbf{Purge committed generated junk}: build artifacts and
  data files committed at the repo root instead of under \texttt{data/}.
\item \textbf{Fix \texttt{.gitignore}}: currently wraps the real ignore
  rules in a literal PowerShell heredoc command that got committed instead
  of its output.
\item \textbf{Delete disk-only bloat} (not committed, safe to remove
  locally): environment archives, unrelated test data, \texttt{\_\_pycache\_\_/}.
\item \textbf{Consolidate duplicated helpers} into \texttt{src/utils/}
  (mostly done -- see \texttt{get\_available\_years},
  \texttt{get\_cached\_shapefile}).
\item \textbf{Resolve naming/content mismatches}: \texttt{display.py},
  \texttt{sealevel2.py}, stale internal docstrings.
\item \textbf{Decide the fate of superseded scripts}:
  \texttt{graficas\_altaresolucion.py}, \texttt{docs/data\_dictionary.md},
  root \texttt{articles/article1.tex}.
\item \textbf{Fill or drop the empty doc stubs} in \texttt{docs/}.
\item \textbf{Populate or drop} the empty
  \texttt{data/processed/anomalias\_\{bogota,medellin,san\_andres\_providencia\}/}
  directories and the unused \texttt{data/indices/}, \texttt{data/zones/}
  locations.
\end{enumerate}

Confirm scope with a maintainer before acting on any of these -- several
things that look like dead scratch code are load-bearing work that simply
never got committed (\Cref{sec:git-gap}).

\end{document}
